% GENERATED FILE. Edit canonical lessons, then run npm run generate:books.
% canonical-source-hash: fnv1a64:a4693955b983c6f6
% canonical-lessons: SA-C496-mangalasutram, SA-C496-vadhuh, SA-C496-parinayati, SA-C496-mandapah, SA-C496-saptapadi

\chapter{A Wedding}
\label{ch:sa-a-wedding}
\hlchaptermodality{496}

\begin{chapteropening}
\textbf{By the end of this chapter:} \emph{I can describe a wedding.}

Complete the last lesson of chapter 496: I can describe a wedding.
\end{chapteropening}

\section[\texorpdfstring{\sk{मंगलसूत्रम्} (maṅgalasūtram)}{मंगलसूत्रम् (maṅgalasūtram)}]{\sk{मंगलसूत्रम्} (maṅgalasūtram) — the marriage thread}

\label{lesson:SA-C496-mangalasutram}



\begin{quote}
Before the new one: say the Sanskrit for holds back, then the Sanskrit for a celibate student.
\end{quote}

\subsection*{You'll want to know: \sk{मंगलसूत्रम्}}
\textbf{\sk{मंगलसूत्रम्}} — \emph{maṅgalasūtram} — \textquotedblleft{}the marriage thread\textquotedblright{}.

\begin{quote}
\sk{विवाहे} \sk{वरः} \sk{मंगलसूत्रं} \sk{बध्नाति।}
\end{quote}

\begin{quote}
\emph{vivāhe varaḥ maṅgalasūtraṃ badhnāti.}
\end{quote}

\begin{quote}
At the wedding the groom ties the marriage thread.
\end{quote}

A neuter noun.

\begin{cousinweb}
\textbf{\sk{मंगलसूत्रम्}} (\emph{maṅgalasūtram}) is built from \emph{maṅgala}, \textquotedblleft{}good luck, a blessing\textquotedblright{}, and \textbf{\sk{सूत्रम्}}, \textquotedblleft{}a thread\textquotedblright{}: MW gives \textquotedblleft{}lucky thread\textquotedblright{}, the marriage thread tied by the bridegroom round the bride's neck (MW).
\end{cousinweb}

\begin{grammarlens}[title={What you've built}]
One more word for describing what you lived and felt.
\end{grammarlens}

\subsection*{Your turn}
\begin{itemize}
\raggedright
  \item \emph{Say it:} \emph{maṅgalasūtram}
  \item \emph{Say it:} \emph{maṅgalasūtram}, once more
  \item \emph{Say it:} \emph{vivāhe varaḥ maṅgalasūtraṃ badhnāti.} — At the wedding the groom ties the marriage thread.
  \item \emph{Recall:} say the Sanskrit for holds back, then the Sanskrit for a celibate student, then say \emph{maṅgalasūtram} again
\end{itemize}

\begin{tcolorbox}[breakable,colback=teal!4,colframe=teal!35!black,title={Before you move on}]
What does \sk{मंगलसूत्रम्} mean? (\textquotedblleft{}The marriage thread\textquotedblright{}.) Say it once more.
\end{tcolorbox}
\section[\texorpdfstring{\sk{वधूः} (vadhūḥ)}{वधूः (vadhūḥ)}]{\sk{वधूः} (vadhūḥ) — a bride}

\label{lesson:SA-C496-vadhuh}



\begin{quote}
Before the new one: say the Sanskrit for a celibate student, then the Sanskrit for the marriage thread.
\end{quote}

\subsection*{You'll want to know: \sk{वधूः}}
\textbf{\sk{वधूः}} — \emph{vadhūḥ} — \textquotedblleft{}a bride\textquotedblright{}.

\begin{quote}
\sk{वधूः} \sk{रक्तां} \sk{शाटिकां} \sk{धारयति।}
\end{quote}

\begin{quote}
\emph{vadhūḥ raktāṃ śāṭikāṃ dhārayati.}
\end{quote}

\begin{quote}
The bride wears a red sari.
\end{quote}

A feminine noun.

\begin{cousinweb}
\textbf{\sk{वधूः}} (\emph{vadhūḥ}) is traced to √vadh, a form of the root √vah, \textquotedblleft{}to carry\textquotedblright{}, by MW: the one led home; MW gives \textquotedblleft{}a bride, a newly married woman\textquotedblright{} (MW).
\end{cousinweb}

\begin{grammarlens}[title={What you've built}]
One more word for describing what you lived and felt.
\end{grammarlens}

\subsection*{Your turn}
\begin{itemize}
\raggedright
  \item \emph{Say it:} \emph{vadhūḥ}
  \item \emph{Say it:} \emph{vadhūḥ}, once more
  \item \emph{Say it:} \emph{vadhūḥ raktāṃ śāṭikāṃ dhārayati.} — The bride wears a red sari.
  \item \emph{Recall:} say the Sanskrit for a celibate student, then the Sanskrit for the marriage thread, then say \emph{vadhūḥ} again
\end{itemize}

\begin{tcolorbox}[breakable,colback=teal!4,colframe=teal!35!black,title={Before you move on}]
What does \sk{वधूः} mean? (\textquotedblleft{}A bride\textquotedblright{}.) Say it once more.
\end{tcolorbox}
\section[\texorpdfstring{\sk{परिणयति} (pariṇayati)}{परिणयति (pariṇayati)}]{\sk{परिणयति} (pariṇayati) — marries}

\label{lesson:SA-C496-parinayati}



\begin{quote}
Before the new one: say the Sanskrit for the marriage thread, then the Sanskrit for a bride.
\end{quote}

\subsection*{You'll want to know: \sk{परिणयति}}
\textbf{\sk{परिणयति}} — \emph{pariṇayati} — \textquotedblleft{}marries\textquotedblright{}.

\begin{quote}
\sk{युवकः} \sk{युवतिं} \sk{परिणयति।}
\end{quote}

\begin{quote}
\emph{yuvakaḥ yuvatiṃ pariṇayati.}
\end{quote}

\begin{quote}
The young man marries the young woman.
\end{quote}

Past (he …): \sk{सः} \sk{युवतिं} \sk{परिणीतवान्}, \emph{saḥ yuvatiṃ pariṇītavān}.

\begin{cousinweb}
\textbf{\sk{परिणयति}} (\emph{pariṇayati}) is built from \emph{pari}, \textquotedblleft{}around\textquotedblright{}, and the root √nī, \textquotedblleft{}to lead\textquotedblright{}: to lead the bride round the fire; MW's ppp \emph{pariṇīta} is \textquotedblleft{}married\textquotedblright{} (MW).
\end{cousinweb}

\begin{grammarlens}[title={What you've built}]
One more word for describing what you lived and felt.
\end{grammarlens}

\subsection*{Your turn}
\begin{itemize}
\raggedright
  \item \emph{Say it:} \emph{pariṇayati}
  \item \emph{Say it:} \emph{pariṇayati}, once more
  \item \emph{Say it:} \emph{yuvakaḥ yuvatiṃ pariṇayati.} — The young man marries the young woman.
  \item \emph{Recall:} say the Sanskrit for the marriage thread, then the Sanskrit for a bride, then say \emph{pariṇayati} again
\end{itemize}

\begin{tcolorbox}[breakable,colback=teal!4,colframe=teal!35!black,title={Before you move on}]
What does \sk{परिणयति} mean? (\textquotedblleft{}Marries\textquotedblright{}.) Say it once more.
\end{tcolorbox}
\section[\texorpdfstring{\sk{मण्डपः} (maṇḍapaḥ)}{मण्डपः (maṇḍapaḥ)}]{\sk{मण्डपः} (maṇḍapaḥ) — a pavilion}

\label{lesson:SA-C496-mandapah}



\begin{quote}
Before the new one: say the Sanskrit for a bride, then the Sanskrit for marries.
\end{quote}

\subsection*{You'll want to know: \sk{मण्डपः}}
\textbf{\sk{मण्डपः}} — \emph{maṇḍapaḥ} — \textquotedblleft{}a pavilion\textquotedblright{}.

\begin{quote}
\sk{वधूः} \sk{मण्डपं} \sk{प्रविशति।}
\end{quote}

\begin{quote}
\emph{vadhūḥ maṇḍapaṃ praviśati.}
\end{quote}

\begin{quote}
The bride enters the pavilion.
\end{quote}

A masculine noun.

\begin{cousinweb}
\textbf{\sk{मण्डपः}} (\emph{maṇḍapaḥ}) MW divides it \emph{maṇḍa-pa} but gives no further derivation, so its origin is uncertain; it gives \textquotedblleft{}an open hall or temporary shed erected on festive occasions, a pavilion\textquotedblright{} (MW).
\end{cousinweb}

\begin{grammarlens}[title={What you've built}]
One more word for describing what you lived and felt.
\end{grammarlens}

\subsection*{Your turn}
\begin{itemize}
\raggedright
  \item \emph{Say it:} \emph{maṇḍapaḥ}
  \item \emph{Say it:} \emph{maṇḍapaḥ}, once more
  \item \emph{Say it:} \emph{vadhūḥ maṇḍapaṃ praviśati.} — The bride enters the pavilion.
  \item \emph{Recall:} say the Sanskrit for a bride, then the Sanskrit for marries, then say \emph{maṇḍapaḥ} again
\end{itemize}

\begin{tcolorbox}[breakable,colback=teal!4,colframe=teal!35!black,title={Before you move on}]
What does \sk{मण्डपः} mean? (\textquotedblleft{}A pavilion\textquotedblright{}.) Say it once more.
\end{tcolorbox}
\section[\texorpdfstring{\sk{सप्तपदी} (saptapadī)}{सप्तपदी (saptapadī)}]{\sk{सप्तपदी} (saptapadī) — the seven steps}

\label{lesson:SA-C496-saptapadi}



\begin{quote}
Before the new one: say the Sanskrit for marries, then the Sanskrit for a pavilion.
\end{quote}

\subsection*{You'll want to know: \sk{सप्तपदी}}
\textbf{\sk{सप्तपदी}} — \emph{saptapadī} — \textquotedblleft{}the seven steps\textquotedblright{}.

\begin{quote}
\sk{सप्तपद्याः} \sk{अनन्तरं} \sk{विवाहः} \sk{पूर्णः} \sk{भवति।}
\end{quote}

\begin{quote}
\emph{saptapadyāḥ anantaraṃ vivāhaḥ pūrṇaḥ bhavati.}
\end{quote}

\begin{quote}
After the seven steps the wedding is complete.
\end{quote}

A feminine noun.

\begin{cousinweb}
\textbf{\sk{सप्तपदी}} (\emph{saptapadī}) is built from \textbf{\sk{सप्त}}, \textquotedblleft{}seven\textquotedblright{}, and \emph{pada}, \textquotedblleft{}a step\textquotedblright{}; MW gives \textquotedblleft{}the seven steps round the sacred fire at the marriage ceremony\textquotedblright{} (MW).
\end{cousinweb}

\begin{grammarlens}[title={What you've built}]
One more word for describing what you lived and felt.
\end{grammarlens}

\subsection*{Your turn}
\begin{itemize}
\raggedright
  \item \emph{Say it:} \emph{saptapadī}
  \item \emph{Say it:} \emph{saptapadī}, once more
  \item \emph{Say it:} \emph{saptapadyāḥ anantaraṃ vivāhaḥ pūrṇaḥ bhavati.} — After the seven steps the wedding is complete.
  \item \emph{Recall:} say the Sanskrit for marries, then the Sanskrit for a pavilion, then say \emph{saptapadī} again
\end{itemize}

\begin{tcolorbox}[breakable,colback=teal!4,colframe=teal!35!black,title={Before you move on}]
What does \sk{सप्तपदी} mean? (\textquotedblleft{}The seven steps\textquotedblright{}.) Say it once more.
\end{tcolorbox}
